%% Full term-theory and algebra/model equivalence proofs.

\section{Lawvere theories of equational presentations}

\begin{annexitem}{annex:p10.lawvere.1}{thm:p10-lawvere-term-theory}{Term tuples, composition, products, and evaluation}
\begin{proof}
\paragraph{Substitution and the quotient arrows.}
Term substitution is defined by the variable and operation
recursion in Definition~\ref{def:p10-term-algebra}.
Induction on a term proves that substitution of its own variable
tuple leaves it unchanged, and that performing two substitutions
in succession equals substituting their componentwise composite.
For a variable these are the defining clauses; an operation
applies the induction hypotheses to its finitely many arguments.
A nullary operation is unchanged in either calculation.
These are syntactic identities, before quotienting.

Changing a component of the second arrow to an
\(\equiv_E^m\)-equal term does not change its composite class,
by the cross-context substitution condition.
If \(r_i\equiv_E^n r'_i\) for every component of the first
arrow, induction on a term \(t\in T_\Sigma(m)\) gives
\[
                         t[r_i/x_i]\equiv_E^n t[r'_i/x_i].
\]
At a variable this is the corresponding component hypothesis,
and at an operation it is congruence compatibility applied
to the induction hypotheses.  Thus composition is well-defined.
Only a finite tuple of representatives is used for a
particular composite.  No simultaneous choice of representatives
for all arrows is required.
The syntactic substitution identities pass to classes and
prove both identity laws and associativity.

There is exactly one zero-component tuple from \(n\) to \(0\).
If \(T_\Sigma(0)\) is empty, there is no arrow \(0\to m\)
for \(m>0\); this does not obstruct any defined composite.
The identity of \(0\) is the empty tuple.
These observations verify the category also in all empty cases.

\paragraph{Products and the structural functor.}
For \(n>0\), the arrows \(p_i^n=([x_i]):n\to1\) have the
following universal property.  Given \(n\) arrows \(k\to1\),
their component classes form a unique tuple \(k\to n\);
composition with \(p_i^n\) recovers its \(i\)-th component.
For \(n=0\) the unique tuple already proves terminality.
Thus \(n\) is the selected \(n\)-fold power of \(1\), with
\(p_0^1=1_1\).

For a finite-set map \(u:m\to n\), assign the variable-selection
tuple \(j(u^{\mathrm{op}})=([x_{u(i)}])_{i<m}\).
The identity map gives the identity tuple.
For \(v:k\to m\), substituting these tuples gives
\[
       j(v^{\mathrm{op}})\circ j(u^{\mathrm{op}})
                         =([x_{u(v(i))}])_{i<k}
                         =j((u\circ v)^{\mathrm{op}}).
\]
This proves functoriality.
The selected structural projections are exactly the just-proved
power projections, so this functor preserves finite products.
In particular, the two block tuples \(n+m\to n,m\) have the
binary product property: concatenate any two tuples with a
common source, and recover them through the block projections.
The blocks determine the concatenated tuple uniquely, including
when one or both blocks are empty.
All term sets, quotient sets, and the countable family of
hom-sets lie in the fixed universe, so the theory is small.

\paragraph{Evaluation is well-defined.}
Fix \(A\models E\).
For every \(n\), define a relation on \(T_\Sigma(n)\) by
\[
       s\mathrel{R_n}t
             \quad\Longleftrightarrow\quad
       s^A(a)=t^A(a)\text{ for every }a\in A^n.
\]
Each \(R_n\) is an equivalence relation compatible with
the operations.  It contains each axiom in its specified
context because \(A\models E\).
The substitution law for evaluation shows that the family
\((R_n)_n\) satisfies the cross-context substitution condition:
evaluating a substituted term first evaluates its tuple of
substituends to an assignment for the original context.
This law is proved by the same term recursion as
Theorem~\ref{thm:p10-term-universal}, and applies to any existing
assignment, including the unique assignment in an empty context.
Therefore leastness gives \(\equiv_E^n\subseteq R_n\).
Term evaluation is independent of the chosen component
representatives.

The substitution law now makes the displayed \(M_A\) preserve
composition, and evaluation of variables gives its identity
law.  For \(n=0\), \(M_A(0)=A^0\) is a singleton.
The cone \(M_A(n)\to M_A(1)\) sends a tuple \(a\in A^n\) to
its singleton-coordinate tuples \((a_i)\in A^1\).
Its comparison
\[
                 A^n\longrightarrow(A^1)^n,\qquad
                 (a_i)_{i<n}\longmapsto((a_i))_{i<n}
\]
is a bijection, with inverse taking the sole coordinate in
each singleton tuple.  These are product cones also when
\(A\) or \(n\) is empty.  Block concatenation likewise gives
the comparison \(A^{n+m}\cong A^n\times A^m\).
Thus \(M_A\) preserves finite products.

\paragraph{Homomorphisms give natural transformations.}
For a homomorphism \(h:A\to B\), induction on a term gives
\[
                            h(t^A(a))=t^B(h^n(a)).
\]
Variables use the coordinate definition of \(h^n\);
operations use preservation by \(h\), including the nullary
case.  Applying this equation to every component of an arrow
gives
\[
                   h^m M_A(([t_i]))=M_B(([t_i]))h^n.
\]
This is naturality.  At object \(0\), \(h^0\) is the unique
function between the two singleton zero-powers.
Powers preserve identity and composition of functions, so
\(A\mapsto M_A,\ h\mapsto M_h\) itself defines a functor.
\end{proof}
\end{annexitem}

\begin{annexitem}{annex:p10.lawvere.2}{thm:p10-lawvere-models}{Recovering operations and all homomorphisms}
\begin{proof}
Write \(K(A)=M_A\) for the evaluation functor of the preceding
theorem, and put \(L(M)=A_M\) with the proposed operations.
The comparison \(c_k^M\) is bijective because \(M\) preserves
the specified power cone.  The definition therefore gives a
function \(A_M^k\to A_M\) for every basic operation, including
\(k=0\).
If a nullary symbol exists, its function from the singleton
zero-power in particular ensures that \(A_M\) is nonempty.

\paragraph{Term interpretation and the equations.}
For \(t\in T_\Sigma(n)\), let \(\bar t:n\to1\) be its
one-component arrow.  We prove by induction that
\[
                         t^{A_M}\circ c_n^M=M(\bar t).
                                                                  \tag{1}
\]
For \(t=x_i\), the left side is the \(i\)-th coordinate of
\(c_n^M\), namely \(M(p_i^n)\), as required.
Suppose \(t=\sigma(t_0,\ldots,t_{k-1})\), and let
\(r=(\bar t_0,\ldots,\bar t_{k-1}):n\to k\) be the tuple
arrow.  Its projections are \(p_i^k r=\bar t_i\), so
\[
                       c_k^M M(r)
                           =(M(\bar t_i))_{i<k}.
                                                                  \tag{2}
\]
For \(k=0\), both sides are the unique map into the singleton
zero-power, so the equation remains valid.
Using (2), the induction hypotheses, and
\(\bar t=\bar\sigma r\), we obtain
\[
 \begin{aligned}
 t^{A_M}c_n^M
   &=M(\bar\sigma)(c_k^M)^{-1}
                         (t_i^{A_M}c_n^M)_{i<k}\\
   &=M(\bar\sigma)(c_k^M)^{-1}c_k^M M(r)\\
   &=M(\bar\sigma r)=M(\bar t).
 \end{aligned}
\]
This proves (1).
For every \((n,s,t)\in E\), the quotient theory has
\(\bar s=\bar t\).
Equation (1) and surjectivity of \(c_n^M\) give
\(s^{A_M}=t^{A_M}\) on every assignment in \(A_M^n\).
Thus \(A_M\models E\).

\paragraph{Natural transformations are exactly homomorphisms.}
For \(\alpha:M\Rightarrow N\), naturality at \(p_i^n\) gives
\[
                        N(p_i^n)\alpha_n=\alpha_1 M(p_i^n).
\]
Taking the tuple of these equations yields
\[
                         c_n^N\alpha_n=(\alpha_1)^n c_n^M.
                                                                  \tag{3}
\]
For \(n=0\), both sides have the same singleton codomain and
are equal by terminal uniqueness.
Naturality at a basic operation arrow \(\bar\sigma:k\to1\),
together with (3), gives
\[
 \begin{aligned}
 \alpha_1\sigma^{A_M}
   &=\alpha_1 M(\bar\sigma)(c_k^M)^{-1}\\
   &=N(\bar\sigma)\alpha_k(c_k^M)^{-1}\\
   &=N(\bar\sigma)(c_k^N)^{-1}(\alpha_1)^k
    =\sigma^{A_N}(\alpha_1)^k.
 \end{aligned}
\]
Hence \(\alpha_1\) is an algebra homomorphism.
Evaluation at object \(1\) preserves transformation identities
and compositions, so \(L\) is a functor.
Equation (3) also proves that no other component of a
transformation is free to vary once \(\alpha_1\) is fixed.

\paragraph{The natural comparison on models.}
For \(r:n\to m\), write \(r=(\bar t_i)_{i<m}\).
The coordinates of \(c_m^M M(r)\) are \(M(\bar t_i)\).
By (1) these are the coordinates of
\(K L(M)(r)c_n^M\).
Thus
\[
                         c_m^M M(r)=K L(M)(r)c_n^M.
                                                                  \tag{4}
\]
Each \(c_n^M\) is a bijection, so they form a natural
isomorphism \(c^M:M\Rightarrow KL(M)\).
For a transformation \(\alpha:M\Rightarrow N\), equation (3)
is precisely
\[
                         c^N\alpha=KL(\alpha)c^M.
\]
Consequently these isomorphisms are natural in the model as
well as in the objects of its theory.

\paragraph{The natural comparison on algebras.}
The carrier of \(LK(A)\) is \(K(A)(1)=A^1\).
Let \(\eta_A:A\to A^1\) send \(a\) to its singleton tuple
\((a)\).  Its inverse takes the sole coordinate.
For a basic \(k\)-ary operation, the construction of \(L\)
and the explicit comparison for \(K(A)\) show that the
operation on \(A^1\) sends \(((a_i))_{i<k}\) to
\((\sigma^A((a_i)_{i<k}))\).
It therefore commutes with \(\eta_A\), including the
nullary case.  Thus \(\eta_A\) is an algebra isomorphism.
For \(h:A\to B\), the equality
\(h^1\eta_A=\eta_B h\) proves naturality.

The functors \(K,L\), with
\(\eta:1_{\operatorname{Alg}_{\mathrm{all}}}\Rightarrow LK\)
and \(\varepsilon_M=(c^M)^{-1}:KL(M)\Rightarrow M\), give
the asserted equivalence.
They also satisfy both triangles.
At object \(n\), \(K(\eta_A)\) sends
\((a_i)\) to \(((a_i))\), which is exactly \(c_n^{K(A)}\);
hence \(\varepsilon_{K(A)}K(\eta_A)\) is the identity.
For the second triangle, \(\eta_{L(M)}\) sends
\(a\in M(1)\) to \((a)\), and
\(L(\varepsilon_M)=(c_1^M)^{-1}\) takes the sole coordinate.
Their composite is the identity of \(L(M)\).

All maps and identities remain valid on empty carriers:
\(\eta_\varnothing\) is the unique bijection
\(\varnothing\to\varnothing^1\), and each zero-power is
still a singleton.
Finally \(A\ne\varnothing\) if and only if \(K(A)(1)=A^1\)
is nonempty, while \(L(M)\) has carrier \(M(1)\).
The functors and all comparison isomorphisms therefore
restrict to the two specified full subcategories.
\end{proof}
\end{annexitem}
